% !PRINTED 132
\weberpage{132}
%<<CONT

Aus der Stetigkeit der Function $f(x)$ folgt aber dann leicht, dass $f(\xi)=0$ sein muss; denn wäre $|f(\xi)|>0$, so könnten, da die $\alpha^{(\nu)}$ dem Werthe $\xi$ beliebig nahe kommen, die absoluten Werthe $a^{(\nu)}$ von $f(\alpha^{(\nu)})$ nicht unter jeden positiven Werth herunter sinken, was wir doch von ihnen nachgewiesen haben.


\begin{center}
{\large \S.\ 40.}

\vspace{0.5em}
{\large\gesperrt{Stetigkeit der Wurzeln.}}
\end{center}

Wir beschliessen diesen Abschnitt mit dem Beweis des Satzes:
\begin{center}
\gesperrt{Die Wurzeln einer algebraischen Gleichung sind stetige Functionen der Coëfficienten.}
\end{center}

Wir haben zunächst die Bedeutung dieses Satzes zu erklären.

Es sei
\begin{equation}
 f(x)=x^n+a_1x^{n-1}+a_2x^{n-2}+\cdots
\tag{1}
\end{equation}
eine ganze rationale Function von $x$ vom $n$ten Grade. Nach dem, was in den vorangegangenen Paragraphen bewiesen ist, lässt sich $f(x)$ in $n$ lineare Factoren zerlegen, die zum Theil einander gleich sein können. Wir setzen, indem wir gleiche Factoren zusammenfassen und die Wurzeln mit $\alpha,\beta,\gamma,
\ldots$ bezeichnen,
\begin{equation}
 f(x)=(x-\alpha)^a(x-\beta)^b(x-\gamma)^c\cdots,
\tag{2}
\end{equation}
worin $a,b,c,
\ldots$ ganze positive Zahlen sind, deren Summe gleich $n$ ist.

Die Wurzeln $\alpha,\beta,\gamma,
\ldots$ werden sich mit den Coëfficienten $a_1,a_2,a_3,
\ldots$ ändern, auch der Grad ihrer Vielfachheit kann ein anderer werden.

Wir bezeichnen die Aenderungen von $a_1,a_2,
\ldots$ mit $\varepsilon_1,\varepsilon_2,
\ldots$ und setzen
\begin{align}
 \varphi(x)&=\varepsilon_1x^{n-1}+\varepsilon_2x^{n-2}+\cdots,
\tag{3}\\
 f(x)+\varphi(x)&=f_1(x).
\tag{4}
\end{align}
Wir umgeben die Punkte $\alpha,\beta,\gamma,
\ldots$ mit Gebieten von beliebiger Kleinheit, jedoch so, dass diese Gebiete sich gegenseitig ausschliessen, etwa dadurch, dass wir die Punkte $\alpha,\beta,\gamma,
\ldots$ durch Kreisperipherien mit den Radien $\varrho,\varrho',\varrho'',
\ldots$ einschliessen, und bezeichnen diese Gebiete durch $(\varrho),(\varrho'),(\varrho''),\ldots$.
